\talkslide{beam-sft}{Safety SFT}{Timing: provisional}{We built the SFT data iteratively: train a model, find prompts that elicit harmful compliance or over-refusal, and add successful attacks to the next round.}{
\begin{tikzpicture}[remember picture,overlay,
 every node/.style={text=Ink,align=left,inner sep=0pt,font=\normalsize},
 trainingbox/.style={draw=Accent!65,fill=Pale,rounded corners=5pt,line width=1pt,
  text width=4cm,minimum height=1.15cm,align=center,inner sep=.12cm},
 flow/.style={-{Latex[length=2.1mm]},draw=Ink,line width=.85pt}]
\begin{scope}[shift={(current page.south west)}]
 \node[trainingbox] (generate) at (2.8,4.75)
  {Generate with policy\\[.08cm]\detail{(deliberative alignment)}};
 \node[trainingbox] (rewrite) at (8,4.75) {Filter + rewrite\\reasoning traces};
 \node[trainingbox,draw=LegacyGreen!65,fill=LegacyGreen!4] (train) at (13.2,4.75)
  {Safety SFT\\[.08cm]\detail{(policy removed)}};
 \draw[flow] (generate.east)--(rewrite.west);
 \draw[flow] (rewrite.east)--(train.west);
\end{scope}
\end{tikzpicture}

\paperref{deliberative}
}
\speak{Supervised fine-tuning pairs a prompt with a desired response and trains the model to imitate that response. Safety examples are mixed with all the other SFT data, not trained in a separate phase at the end. We built the SFT data iteratively: train a model, find prompts that elicit harmful compliance or over-refusal, and add successful attacks to the next round. We cover harmful and harmless requests, including direct requests and jailbreaks, while training against over-refusal. SFT teaches the model to imitate these desired responses, but other sampled responses can still violate the safety policy. This is why we also use safety tasks in RL to train on the model’s own responses. Following Deliberative Alignment, a model first generated reasoning and answers with the safety policy in context. A judge filtered poor behavior, and another model rewrote the traces to express general safety reasoning. The training examples removed the supplied policy. See \href{https://arxiv.org/abs/2412.16339}{Guan, Joglekar, Wallace et al., Deliberative Alignment (2024)} for policy-conditioned reasoning-trace generation and policy-free SFT. The Beam report describes its additional filtering and rewriting step. Safety examples were 1.6 percent of the alignment teacher's SFT dataset, not 1.6 percent of all Beam training. Source detail: \href{https://reflection.ai/documents/beam-model-tech-report}{Beam technical report}.}
